% GENERATED FILE. Edit canonical lessons, then run npm run generate:books.
% canonical-source-hash: fnv1a64:b4c000c1f24e4ad3
% canonical-lessons: ZH-C216-zaogao, ZH-C216-qiguai, ZH-C216-tou, ZH-C216-jingcha, ZH-C216-baojing

\chapter{Something Has Gone Wrong}
\label{ch:zh-something-has-gone-wrong}
\hlchaptermodality{216}

\begin{chapteropening}
\textbf{By the end of this chapter:} \emph{I can say something has gone wrong or is strange, that someone has stolen something, and call the police.}

Complete the last lesson of chapter 216: I can say something has gone wrong or is strange, that someone has stolen something, and call the police.
\end{chapteropening}

\section[zāogāo]{zāogāo — oh no; what a mess}

\label{lesson:ZH-C216-zaogao}



\begin{quote}
Before the new one: say the Mandarin for laundry shop, then the Mandarin for to dry-clean.
\end{quote}

\subsection*{You'll want to know: zāogāo}
\emph{zāogāo} — \textquotedblleft{}oh no; what a mess\textquotedblright{}.

\begin{quote}
\emph{zāogāo wǒ de qiánbāo diū le} — Oh no, I've lost my wallet!
\end{quote}

Its characters come later; for now it is a word for the ear.

\begin{cousinweb}
\emph{zāogāo} is built from \emph{zāo} 'dregs' and \emph{gāo} 'cake', both inherited from Middle Chinese (Guangyun, in Baxter's transcription): \emph{tsaw}, glossed 'dregs', and probably the \emph{kaw} that the dictionary writes with another form of the character, glossed with a word for gruel. How the pair came to mean a mess is not settled.
\end{cousinweb}

\begin{grammarlens}[title={What you've built}]
One more word for saying no and asking.
\end{grammarlens}

\subsection*{Your turn}
\begin{itemize}
\raggedright
  \item \emph{Say it:} \emph{zāogāo}
  \item \emph{Say it:} \emph{zāogāo}, once more
  \item \emph{Say it:} \emph{zāogāo wǒ de qiánbāo diū le} — Oh no, I've lost my wallet!
  \item \emph{Recall:} say the Mandarin for laundry shop, then the Mandarin for to dry-clean, then say \emph{zāogāo} again
\end{itemize}

\begin{tcolorbox}[breakable,colback=teal!4,colframe=teal!35!black,title={Before you move on}]
What does zāogāo mean? (\textquotedblleft{}Oh no; what a mess\textquotedblright{}.) Say it once more.
\end{tcolorbox}
\section[qíguài]{qíguài — strange; odd}

\label{lesson:ZH-C216-qiguai}



\begin{quote}
Before the new one: say the Mandarin for to dry-clean, then the Mandarin for oh no.
\end{quote}

\subsection*{You'll want to know: qíguài}
\emph{qíguài} — \textquotedblleft{}strange; odd\textquotedblright{}.

\begin{quote}
\emph{zhège rén hěn qíguài} — That person is strange.
\end{quote}

Its characters come later; for now it is a word for the ear.

\begin{cousinweb}
\emph{qíguài} is built from \emph{qí} and \emph{guài}, both inherited from Middle Chinese, \emph{gje} and \emph{kweajH} (Guangyun, in Baxter's transcription), which glosses the first 'different' and the second 'strange'.
\end{cousinweb}

\begin{grammarlens}[title={What you've built}]
One more word for saying no and asking.
\end{grammarlens}

\subsection*{Your turn}
\begin{itemize}
\raggedright
  \item \emph{Say it:} \emph{qíguài}
  \item \emph{Say it:} \emph{qíguài}, once more
  \item \emph{Say it:} \emph{zhège rén hěn qíguài} — That person is strange.
  \item \emph{Recall:} say the Mandarin for to dry-clean, then the Mandarin for oh no, then say \emph{qíguài} again
\end{itemize}

\begin{tcolorbox}[breakable,colback=teal!4,colframe=teal!35!black,title={Before you move on}]
What does qíguài mean? (\textquotedblleft{}Strange; odd\textquotedblright{}.) Say it once more.
\end{tcolorbox}
\section[tōu]{tōu — to steal}

\label{lesson:ZH-C216-tou}



\begin{quote}
Before the new one: say the Mandarin for oh no, then the Mandarin for odd.
\end{quote}

\subsection*{You'll want to know: tōu}
\emph{tōu} — \textquotedblleft{}to steal\textquotedblright{}.

\begin{quote}
\emph{yǒu rén tōule wǒ de zìxíngchē} — Somebody has stolen my bike.
\end{quote}

Its characters come later; for now it is a word for the ear.

\begin{cousinweb}
\emph{tōu} is inherited from Middle Chinese \emph{thuw} (Guangyun, in Baxter's transcription), glossed 'to rob'.
\end{cousinweb}

\begin{grammarlens}[title={What you've built}]
One more word for saying no and asking.
\end{grammarlens}

\subsection*{Your turn}
\begin{itemize}
\raggedright
  \item \emph{Say it:} \emph{tōu}
  \item \emph{Say it:} \emph{tōu}, once more
  \item \emph{Say it:} \emph{yǒu rén tōule wǒ de zìxíngchē} — Somebody has stolen my bike.
  \item \emph{Recall:} say the Mandarin for oh no, then the Mandarin for odd, then say \emph{tōu} again
\end{itemize}

\begin{tcolorbox}[breakable,colback=teal!4,colframe=teal!35!black,title={Before you move on}]
What does tōu mean? (\textquotedblleft{}To steal\textquotedblright{}.) Say it once more.
\end{tcolorbox}
\section[jǐngchá]{jǐngchá — police officer; the police}

\label{lesson:ZH-C216-jingcha}



\begin{quote}
Before the new one: say the Mandarin for odd, then the Mandarin for to steal.
\end{quote}

\subsection*{You'll want to know: jǐngchá}
\emph{jǐngchá} — \textquotedblleft{}police officer; the police\textquotedblright{}.

\begin{quote}
\emph{jǐngchá lái le} — The police are here.
\end{quote}

Its characters come later; for now it is a word for the ear.

\begin{cousinweb}
\emph{jǐngchá} is built from \emph{jǐng} 'to warn' and \emph{chá} 'to watch over', both inherited from Middle Chinese, \emph{kjaengX} and \emph{tsrheat} (Guangyun, in Baxter's transcription), which glosses the first 'to be on guard' and the second 'to supervise'.
\end{cousinweb}

\begin{grammarlens}[title={What you've built}]
One more word for saying no and asking.
\end{grammarlens}

\subsection*{Your turn}
\begin{itemize}
\raggedright
  \item \emph{Say it:} \emph{jǐngchá}
  \item \emph{Say it:} \emph{jǐngchá}, once more
  \item \emph{Say it:} \emph{jǐngchá lái le} — The police are here.
  \item \emph{Recall:} say the Mandarin for odd, then the Mandarin for to steal, then say \emph{jǐngchá} again
\end{itemize}

\begin{tcolorbox}[breakable,colback=teal!4,colframe=teal!35!black,title={Before you move on}]
What does jǐngchá mean? (\textquotedblleft{}Police officer; the police\textquotedblright{}.) Say it once more.
\end{tcolorbox}
\section[bàojǐng]{bàojǐng — to call the police}

\label{lesson:ZH-C216-baojing}



\begin{quote}
Before the new one: say the Mandarin for to steal, then the Mandarin for police officer.
\end{quote}

\subsection*{You'll want to know: bàojǐng}
\emph{bàojǐng} — \textquotedblleft{}to call the police\textquotedblright{}.

\begin{quote}
\emph{wǒmen yào bàojǐng} — We're going to call the police.
\end{quote}

Its characters come later; for now it is a word for the ear.

\begin{cousinweb}
\emph{bàojǐng} is built from \emph{bào} 'to report' and \emph{jǐng} 'to warn', both inherited from Middle Chinese, \emph{pawH} and \emph{kjaengX} (Guangyun, in Baxter's transcription).
\end{cousinweb}

\begin{grammarlens}[title={What you've built}]
One more word for saying no and asking.
\end{grammarlens}

\subsection*{Your turn}
\begin{itemize}
\raggedright
  \item \emph{Say it:} \emph{bàojǐng}
  \item \emph{Say it:} \emph{bàojǐng}, once more
  \item \emph{Say it:} \emph{wǒmen yào bàojǐng} — We're going to call the police.
  \item \emph{Recall:} say the Mandarin for to steal, then the Mandarin for police officer, then say \emph{bàojǐng} again
\end{itemize}

\begin{tcolorbox}[breakable,colback=teal!4,colframe=teal!35!black,title={Before you move on}]
What does bàojǐng mean? (\textquotedblleft{}To call the police\textquotedblright{}.) Say it once more.
\end{tcolorbox}
